% Continuing ninth-hour work: exact finite involution fusion.
% Root owns A.4.56. This fragment is retained separately until audited.
\subsection{A finite fusion refinement for the named elliptic involution}

Keep the actual normalized configuration, labelled involution and
ordinary groups $X_n,x_n,K_n,\mathcal M_n$ of A.4.54, for
$9\leq n\leq12$.  Let
$q_0:\mathcal M_n\twoheadrightarrow Q_0$ be a specified finite
ordinary quotient.  Apply the construction (ICR.4)--(ICR.15)
first to $q_0$, and then to its resulting actual quotient $\eta_0$.
This gives actual homomorphisms and projections
\[
 \mathcal M_n\xrightarrow{\eta_1}E_1
       \xrightarrow{p_{10}}E_0\xrightarrow{p_0}Q_0,
 \qquad \eta_0=p_{10}\eta_1,\qquad q_0=p_0\eta_0.
                                                               \tag{ICF.1}
\]
Write $V_0=\ker(\eta_0|_{K_n})$ and
$V_1=\ker(\eta_1|_{K_n})$.  The second construction starts with
$U_0=V_0$: its added label and relative-homology coordinates
already factor through $E_0$.  Its $b=2n-1$ successive
characteristic subgroups end at $U_b=V_1$.

Restrict $\eta_1$ to the actual linear section
$s\Gamma(6)$.  Induce this homomorphism through the $144$
cosets of $\Gamma(6)$ in $\operatorname{SL}_2(\mathbb Z)$.
Denote the induced finite period homomorphism by $\Omega_1$.
Its permutation part retains reduction modulo six and its
identity coordinate on $\Gamma(6)$ is the original
$\eta_1\circ s$.  Form the actual image
\[
 \eta_1^*=(\eta_1,\Omega_1\circ\operatorname{per}),
 \qquad E_1^*=\operatorname{im}\eta_1^*.
                                                               \tag{ICF.2}
\]
The projection $\pi_{10}:E_1^*\to E_0$ is induced by (ICF.1).
All these are homomorphisms of finite ordinary quotient groups;
their source homomorphisms also extend continuously to the full
profinite completion of the displayed ordinary mapping class group.

\begin{result}{A.4.56}{Proved}{An explicit involution refinement fuses every finite complement at the ordinary centralizer output}
The subgroup with identity label and identity $\Omega_1$ period
in $E_1^*$ is exactly $\eta_1^*(K_n)=K_n/V_1$.
Consequently the fibre with label $\sigma$ and period
$\Omega_1(-I)$ is the exact coset
\[
                         (K_n/V_1)\,\eta_1^*(\iota).
                                                               \tag{ICF.3}
\]
Every element $v$ of order two in (ICF.3) has an actual pure,
orientation-preserving ordinary mapping class
$\alpha\in P\Gamma_{1,n}$ such that
\[
                \pi_{10}(v)=\eta_0(\alpha\iota\alpha^{-1}).
                                                               \tag{ICF.4}
\]
The construction uses $2n-1$ cohomology-killing covers after
the centralizer-recognition output.  It retains the ordered
two-torsion assignments and the actual deck displacement of a
fixed point in the lower cover.  Its size satisfies
\[
                         |E_1^*|\leq144\,|E_1|^{145},
                                                               \tag{ICF.5}
\]
with $|E_1|$ bounded by two successive applications of the
explicit function (ICR.4),(ICR.7), starting at $|Q_0|$.
\end{result}

\begin{proof}
\par\medskip\noindent\textbf{The exact period fibre.}\quad
The induction kernel for $\Omega_1$ is contained in $\Gamma(6)$.
For $A$ in that kernel, its identity coordinate also gives
$\eta_1(s(A))=1$.  Take an ordinary representative of an
element with trivial label and trivial $\Omega_1$ period.
Its period $A$ therefore belongs to $\Gamma(6)$, and its
ordinary decomposition in (ICR.12) is $k s(A)$ for $k\in K_n$.
Both coordinates of $\eta_1^*(s(A))$ are trivial, so its image
is precisely $\eta_1^*(k)$.  Conversely every $k\in K_n$ has
the required label and period.  The kernel on $K_n$ is exactly
$V_1$, proving the subgroup assertion.  Multiplication by the
actual involution proves (ICF.3), with its specified order-two
section.  The induced period image has order at most
$144|E_1|^{144}$, proving (ICF.5).

\par\medskip\noindent\textbf{The lower fixed point.}\quad
Let $Y_i$ be the finite equivariant cell model of the cover of
$X_n$ corresponding to $U_i$, from (ICR.4),(ICR.10).
It has dimension at most $d=2n-2$ and is connected.
The element $v\in(K_n/V_1)\rtimes\langle\iota\rangle$
acts as an involution on $Y_b$ and, by projection, on every
$Y_i$.  The covering maps are equivariant for this particular
involution.  They kill all positive ordinary mod-two cohomology,
by (ICR.10), independently of the chosen complement.

Retain the actual cochain complex
\[
       D(Y_i)=\operatorname{Cone}
                   \bigl(\mathbb F_2\longrightarrow
                              C^*(Y_i;\mathbb F_2)\bigr),
                                                               \tag{ICF.6}
\]
with its constant term in degree minus one.  Connectedness makes
it equivariantly quasi-isomorphic to
$C^*(Y_i;\mathbb F_2)/\mathbb F_2$, supported in degrees
$0,\ldots,d$, and its cohomology is the positive ordinary
cohomology of $Y_i$.  The complete $C_2$-resolution applied to
(ICF.6) gives the same bounded-filtration Tate comparison as
in (ICR.11).  Each covering map is zero on the ordinary-cohomology
$E_2$ page.  Composing $d+1=b$ such maps is therefore zero on
every Tate abutment group.

Restriction to the fixed subcomplex is a Tate isomorphism:
every complementary cell lies in a free $C_2$-orbit, whose
complete-resolution complex is acyclic.  The cone retains its
constant term under this restriction.  On the fixed subcomplex
the horizontal differentials are zero, so the resulting map on
the summand of degree minus one is zero.  The lower fixed locus
$Y_0^v$ is nonempty.  Indeed emptiness would make every $Y_i^v$
empty by equivariance.  In that case all complexes
$D(Y_i^v)$ would consist of the constant $\mathbb F_2$ in degree
minus one and all the maps would be the identity.  Their degree
minus-one map is nonzero, contradicting the proved vanishing.
This proves existence in the lower finite cover without assuming
any fixed point in the upper one.

\par\medskip\noindent\textbf{Every ordered branch assignment is positively realized.}\quad
Choose a fixed point $y\in Y_0^v$ and let $\xi$ be its image
in $X_n^\iota$.  Its $k-1$ fixed-label coordinates are an ordered
list of distinct nonzero points of $T[2]$, and its other eight
coordinates form four ordered pairs, by (ICR.8)--(ICR.9).
The reduction of $\Gamma(3)$ modulo two is all
$\operatorname{SL}_2(\mathbb F_2)$: the actual integral matrices
\[
 U^3=\begin{pmatrix}1&3\\0&1\end{pmatrix},\qquad
 L^3=\begin{pmatrix}1&0\\3&1\end{pmatrix}
                                                               \tag{ICF.7}
\]
belong to $\Gamma(3)$ and their reductions generate the six
matrices of that group.  Its action on the three nonzero vectors
is the full permutation action.  It is transitive on the ordered
lists of length $0,1,2,3$ of distinct such vectors.  Select an
actual word $A\in\langle U^3,L^3\rangle$ carrying the reference
fixed-label assignment to that of $\xi$.

The linear map $\ell_A:T\to T$ is positive, fixes every
original three-division point individually, commutes with
$z\mapsto-z$, and has precisely the required fixed-label images.
On the sphere $T/\{\pm1\}$, move the four images of its paired
orbits to the four paired orbits of $\xi$ by an isotopy fixing
the four branch points.  The four independent branch monodromies
in (ICR.9) prescribe the four lift signs.  Its identity-started
lift $\chi$ to the torus commutes with $\iota$, fixes every
point of $T[2]$, and has compact period identity.
Thus $\ell=\chi\ell_A$ is a positive torus homeomorphism,
commutes with the actual linear involution, and carries every
reference labelled point to its corresponding coordinate of
$\xi$.

Choose a path $\gamma:x_n\to\xi$ in the normalized configuration
space and extend it to a torus isotopy fixing the origin.  Let
$\psi$ be its endpoint.  It has compact period identity and
carries the same labelled reference configuration to $\xi$.
Consequently the actual mapping class
\[
                         \alpha_\xi=\psi^{-1}\ell
                                     \in P\Gamma_{1,n}
                                                               \tag{ICF.8}
\]
is pure and positive, with period $A$.  The path model for the
evaluation fibration identifies its conjugate section with the
literal crossed word
\[
                [\gamma*(\iota\gamma)^{-1}]\,\iota
                    =\alpha_\xi\iota\alpha_\xi^{-1}
                   \quad\hbox{in }K_n\rtimes\langle\iota\rangle.
                                                               \tag{ICF.9}
\]
Here (ICF.9) is the actual isotopy comparison of (CFF.16), with
the named involution in place of the named translation.  The two
evaluation paths are exactly $\gamma$ and $\iota\gamma$; their
concatenation is a loop at $x_n$.  The commuting map $\ell$
identifies the terminal involution and contributes no additional
path.  This also fixes the conjugation direction in (ICF.9).

The chosen point $y$ may differ from the endpoint of the lift
of $\gamma$ by a deck element $\bar k\in K_n/V_0$.
Take an ordinary lift $k\in K_n$ of that element.
The conjugate of (ICF.9) by $\bar k$ and the projection of $v$
are two lifts of the same involution to the same connected
cover, both fixing $y$.  Their quotient is a deck transformation
fixing a point and hence is the identity.  Therefore they are
equal on $Y_0$.  Taking $\alpha=k\alpha_\xi$ proves (ICF.4)
through the exact identification $K_n/V_0=\eta_0(K_n)$.
Both factors are pure and positive, so all asserted markings
and the orientation are retained.
\end{proof}

The result is a statement about the specified finite quotients
and all order-two elements in their exact fibre (ICF.3).
Applying it to a transported tower action also requires that
action's actual label and finite-period coordinates.  Those
coordinates are separate source maps and are not inferred from
an assumed stability of an arbitrary quotient kernel.
